\chapter{课程设计流程}\label{s:process}

大多数人设计课程是这样的：

\begin{enumerate}

\item
  有人请你教一样你几乎不懂、或者好几年没想过的东西。

\item
  你开始写幻灯片，讲解你对这个学科的所知。

\item
  两三周后，
  你根据到目前为止讲过的内容编一份作业。

\item
  第 3 步重复好几次。

\item
  你熬到凌晨，编出一份期末考试，
  并向自己保证下次会更有条理。

\end{enumerate}

一种更有效的办法，在精神上类似于一种叫
\gref{g:test-driven-development}{测试驱动开发}（TDD）的编程实践。
用 TDD 的程序员不会先写软件、再测试它能正确工作。
相反，
他们先写测试，
然后只写刚好能让这些测试通过的新软件。

TDD 之所以有效，是因为写测试迫使程序员精确地说清
自己想达成什么、以及「完成」是什么样子。
TDD 还能防止无休止的打磨：
测试一通过，你就停止写代码。
最后，
它降低了确认偏差的风险：
一个还没写这份软件的人，
会比一个刚投入了好几小时苦干、心心念念就想收工的人更客观。

一种类似的、叫做\gref{g:backward-design}{逆向设计}的方法，用在课程设计上效果很好。
这个方法由~\cite{Wigg2005,Bigg2011,Fink2013}各自独立提出，
并在~\cite{McTi2013}里做了总结。
简化之后，它的步骤是：

\begin{enumerate}

\item
  创建或复用学习者画像（下一节会讨论），
  弄清你想帮的是谁、什么会吸引他们。

\item
  头脑风暴，大致想清楚你想覆盖什么、
  打算怎么做、
  预期会遇到哪些问题或误解、
  哪些东西\emph{不}会包括进去，
  等等。
  在这个阶段，画概念图能帮上大忙（\secref{s:memory-concept-maps}）。\index{概念图}

\item
  创建一份总结性评价（\secref{s:models-formative-assessment}）\index{总结性评价}
  来界定你的总体目标。
  它可以是一门课的期末考试，
  也可以是一天工作坊的结业项目；
  不管它的形式和规模如何，
  它都能比你列成要点的目标清单更清楚地显示你希望达到哪里。

\item
  创建形成性评价，\index{形成性评价}
  让学习者有机会练习他们正在学的东西。
  它们也会告诉你（以及他们）是否在取得进展、
  以及需要把注意力放在哪里。
  最好的做法，是把上一步里那份总结性评价用到的知识和技能逐条列出，
  然后为每一条至少创建一个形成性评价。

\item
  根据这些形成性评价的复杂度、
  它们之间的依赖关系、
  以及各主题能在多大程度上激发学习者的动机（\secref{s:motivation-authentic}），
  给它们排序，从而形成课程大纲。

\item
  编写材料，把学习者从一个形成性评价带到下一个。
  每一小时的授课应该由三到五个这样的片段组成。

\item
  写一份课程的概要描述，
  帮目标受众找到它、
  并判断它是否适合自己。

\end{enumerate}

这个方法有助于让教学始终聚焦在目标上。
它也确保学习者在课程结束时不会遇到
他们没准备好的东西。

\begin{aside}{扭曲的激励}
  逆向设计\emph{不}等同于\gref{g:teaching-to-the-test}{应试教学}。
  使用逆向设计时，
  老师设定目标是为了辅助课程设计；
  他们可能永远不会真正考那份自己写的期末考试。
  而在许多学校体系里，
  恰恰相反，
  一个外部权威为所有学习者规定评价标准，
  不管他们各自的情况如何。
  这些总结性评价的结果直接关系到老师的薪水和晋升，
  这就意味着老师有动力去盯着让学习者通过考试，
  而不是帮他们学习。

  \cite{Gree2014}主张，把重点放在考试和测量上，
  对那些有权制定考试的人很有吸引力，
  但除非同时给老师提供支持，让他们根据考试结果做出改进，
  否则不大可能改善结果。
  后者往往缺失，因为
  大型组织通常看重整齐划一胜过产出效率~\cite{Scot1998}。
\end{aside}

逆向设计被描述成一个序列，
但它几乎从不按那种方式来做。
例如，
我们可能在写一道 MCQ 时忽然想到点什么，
于是改变想教的内容；
或者在有了课程大纲之后，重新评估我们想帮的是谁。
不过，
我们留下的笔记应该按上面描述的顺序来呈现，
好让在我们之后使用或维护这门课的人能重走我们的思路
\cite{Parn1986}。

\seclbl{学习者画像}{s:process-personas}

逆向设计流程的第一步，是弄清你的受众是谁。
一种办法是写两三个
\grefdex{g:learner-persona}{学习者画像}{学习者画像}，
就像\secref{s:intro-audience}里那些。
这项技术是从用户体验设计师那里借来的，
他们会为典型用户创建简短档案，
以帮助自己思考受众。

一个学习者画像包括：

\begin{enumerate}

\item
  此人的大致背景；

\item
  他已经会什么；

\item
  他想做什么；
  以及

\item
  他有什么特殊需求。

\end{enumerate}

\secref{s:intro-audience}里的画像具备上面这四点，
外加一段关于本书将如何帮到他们的简短说明。
一个办周末 Python 工作坊的志愿者小组，其学习者画像可能是：

\begin{enumerate}

\item
  Jorge 刚从哥斯达黎加搬到加拿大，学习农业工程。
  他加入了学院足球队，
  正盼着学打冰球。

\item
  除了用 Excel、Word 和上网，
  Jorge 此前最重要的计算机经验，
  是帮他姐姐为家里的生意建一个 WordPress 网站。

\item
  Jorge 想用一台手持设备测量附近农场的土壤性质，
  把数据发到自己的电脑上。
  现在他得用 Excel 打开每个数据文件、
  删掉第一列和最后一列、
  再对剩下的部分算一些统计量。
  他接下来几个月至少得采集 600 次测量，
  实在不想每一个都手工做这些步骤。

\item
  Jorge 英语阅读没问题，
  但有时跟不上满是行话的口头交谈。

\end{enumerate}

与其为每一节课或每一门课都写新画像，
老师通常会创建并共享六七个
覆盖他们可能教到的所有人的画像，
然后从中挑几个来描述特定材料的受众。
这样使用的画像，成了讨论设计问题时一种方便的简写：
互相交流时，老师们可以说：
「Jorge 会明白我们为什么这么做吗？」
或者，
「Jorge 会遇到哪些安装问题？」

\begin{aside}{是他们的目标，不是你的}
  画像描述的应该始终是学习者想做什么，
  而不是你认为他们实际需要什么。
  问问你自己，他们在网上搜索什么；
  那多半不会包含他们还不懂的行话，
  所以作为教学设计者，你要做的事之一，
  就是想办法让你的课程能被搜到。\index{可发现性!课程}
\end{aside}

\seclbl{学习目标}{s:process-objectives}

形成性评价和总结性评价帮助老师弄清他们要教什么，
但为了把这一点传达给学习者和其他老师，
一份课程描述还应该有
\grefdex{g:learning-objective}{学习目标}{学习目标}。
它们有助于确保
每个人对一节课想要达成什么有一致的理解。
例如，
「理解 Git」这样一句话，可能意味着下面任何一种意思：

\begin{itemize}

\item
  学习者能说出像 Git 这样的版本控制系统
  在哪些三个方面优于 Dropbox 这样的文件共享工具，
  以及在哪两个方面不如它们。

\item
  学习者能用桌面 GUI 工具
  把改动过的文件提交到一个 Git 仓库。

\item
  学习者能解释什么是分离头指针（detached HEAD），
  并用命令行操作从这种状态中恢复。

\end{itemize}

\begin{aside}{目标与成效}
  学习目标是一节课努力要达成的东西。
  \gref{g:learning-outcome}{学习成效}是它实际达成的东西，
  也就是学习者真正带走的东西。
  因此，总结性评价的作用就是
  把学习成效与学习目标作比较。
\end{aside}

一个学习目标描述的是：学习者成功完成一节课后，
将如何展示自己学到的东西。
更具体地说，
它有一个\emph{可测量或可验证的动词}，说明学习者将做什么，
并规定\emph{可接受表现的标准}。
写这些起初可能显得束手束脚，
但从长远看，它们会让你、
你的同行
和你的学习者都更满意：
你最终会得到清晰的教学和评估准则，
而你的学习者会因为有了明确的预期而感激你。

理解什么才是一个好的学习目标，一个办法
是看看一个差的目标如何被改进：

\begin{itemize}

\item
  \emph{学习者将获得学习良好编程实践的机会。}\\
  这描述的是一节课的内容，
  而不是成功学习者的特征。\\

\item
  \emph{学习者将对良好编程实践有更好的体会。}\\
  这既没有以主动动词开头，也没有界定学习的层次，
  而且学习的对象没有语境、也不具体。\\

\item
  \emph{学习者将理解如何用 R 编程。}\\
  虽然它以主动动词开头，
  但没有界定学习的层次，
  而且学习的对象仍然太含糊，没法评估。\\

\item
  \emph{学习者将用 R 写出单页的数据分析脚本，来读取、筛选和汇总表格数据。}\\
  它以主动动词开头，
  界定了学习的层次，
  并提供了语境，以确保成效可以被评估。

\end{itemize}

说到挑选动词，
很多老师使用\gref{g:blooms-taxonomy}{布鲁姆分类法}。
它最初发表于 1956 年，并在世纪之交做了更新~\cite{Ande2001}，
是一个被广泛使用的、讨论理解层次的框架。
它最新的形式有六个类别；
下面列出的是通常用于为每个类别所写学习目标的一些动词：

\begin{description}

\item[记忆：]
  通过回忆事实、术语、基本概念和答案，展现对先前所学材料的记忆。
  \emph{（认出、列出、描述、命名、找出）}

\item[理解：]
  通过组织、比较、翻译、解释、给出描述和陈述中心思想，
  展现对事实和想法的理解。
  \emph{（解释、总结、复述、分类、说明）}

\item[应用：]
  通过以不同的方式运用已获得的知识、事实、技术和规则，解决新问题。
  \emph{（构建、识别、使用、规划、选择）}

\item[分析：]
  通过识别动机或原因，考查并把信息拆解成部分；
  作出推断，并找到支持概括的证据。
  \emph{（比较、对比、简化）}

\item[评价：]
  通过依据一套标准对信息、想法的有效性或作品的质量作出判断，
  来提出并捍卫观点。
  \emph{（核对、选择、批评、证明、评级）}

\item[创造：]
  通过以新的模式组合元素或提出替代方案，把信息以不同的方式汇集起来。
  \emph{（设计、构造、改进、改编、最大化、解决）}

\end{description}

布鲁姆分类法几乎出现在每一本教育教科书里，
但~\cite{Masa2018}发现，
即便是经验丰富的教育者，也很难就如何对具体事物分类达成一致。
不过，那些动词仍然有用，
一步一步构建理解这个观念也一样有用：
正如 Daniel Willingham 所说，\index{Willingham, Daniel}
人没有可思考的东西就没法思考~\cite{Will2010}，
而这个分类法能帮老师确保学习者在需要的时候，手头有那些可思考的东西。

思考学习目标的另一种方式来自~\cite{Fink2013}，
它从学习要在学习者身上产生什么改变来定义学习。
\gref{g:finks-taxonomy}{芬克分类法}也有六个类别，
但与布鲁姆不同，它们是互补的而不是层级式的：

\begin{description}

\item[基础知识：]
  理解和记住信息与想法。
  \emph{（记住、理解、识别）}

\item[应用：]
  技能、批判性思维、项目管理。
  \emph{（使用、解决、计算、创造）}

\item[整合：]
  把想法、学习经验和真实生活联系起来。
  \emph{（连接、关联、比较）}

\item[人的维度：]
  了解自己和他人。
  \emph{（开始把自己看成、从……角度理解他人、决定成为）}

\item[关怀：]
  发展新的感受、兴趣和价值观。
  \emph{（对……感到兴奋、准备好去、看重）}

\item[学会学习：]
  成为更好的学习者。
  \emph{（识别……的信息来源、提出关于……的有用问题）}

\end{description}

为一门 HTML 和 CSS 入门课程，
基于这个分类法写出的一组学习目标可能是：

\begin{itemize}

\item
  解释什么是 CSS 属性，以及 CSS 选择器如何工作。

\item
  用常见标签和 CSS 属性为网页设置样式。

\item
  比较并对比用 HTML 和 CSS 写作
  与用桌面出版工具写作。

\item
  识别并纠正示例网页中
  会让视障人士难以与之交互的问题。

\item
  描述你最喜欢的网站中特别吸引你的设计特点，
  并解释为什么。

\item
  描述你最常用的两个关于 CSS 的在线信息来源，
  并解释你喜欢它们什么。

\end{itemize}

\seclbl{可维护性}{s:process-maintainability}

一门课创建出来之后，总得有人维护它，
而如果它以可维护的方式搭建起来，维护起来会容易得多。
但「可维护」到底是什么意思？\index{可维护性（课程的）}
简短的回答是：如果更新一门课比替换它更便宜，它就是可维护的。
这个等式取决于四个因素：

\begin{description}

\item[课程设计被记录得有多好。]
  如果做维护的人不知道（或记不得）
  这门课想达成什么、
  或者为什么按这个特定顺序引入主题，
  那么更新它就会花更多时间。
  使用逆向设计的理由之一，
  就是记录下「每门课为什么是现在这样」的各项决策。

\item[合作者在技术上协作有多容易。]
  老师通常靠互相邮寄 PowerPoint 文件，
  或者把它们放到共享盘里来共享材料。
  像 \hreffoot{http://docs.google.com}{Google Docs} 这样的协作文档工具和 wiki
  是一大进步，
  因为它们允许多人更新同一份文档、并对别人的更新发表评论。
  程序员使用的版本控制系统，
  比如 \hreffoot{http://github.com}{GitHub}，
  是另一种做法。
  它们让任意多的人独立工作，
  然后以可控、可审阅的方式合并各自的改动。
  遗憾的是，
  版本控制系统学习曲线很陡，
  而且处理不了常见的办公文档格式。

\item[人们有多愿意协作。]
  构建一个课程版维基百科所需的工具已经存在二十年了，
  但大多数老师仍然不会像写和分享百科条目那样
  去写和分享课程。

\item[分享实际上有多有用。]
  \gref{g:reusability-paradox}{复用悖论}指出，
  一个学习对象越可复用，
  它的教学效果就越差~\cite{Wile2002}。
  原因是，一门好课更像一部小说，而不是一个程序：
  它的各个部分紧紧耦合，而不是彼此独立黑箱。
  因此直接复用也许是错误的目标；
  如果我们设法让课程更容易重新混编，也许能走得更远。

\end{description}

如果复用悖论是对的，
那么当被协作的东西很小时，协作就更可能发生。
这与 Mike Caulfield 的\index{Caulfield, Mike}
\hreffoot{https://hapgood.us/2016/05/13/choral-explanations/}{合唱式解释}理论很契合，
该理论主张，像 \hreffoot{https://stackoverflow.com/}{Stack Overflow} 这样的网站之所以成功，
是因为它们为每个问题提供了一合唱团式的答案，
每一个都最适合略有不同的提问者。
如果这是对的，
那么未来的课程也许就是由社区整理的问答库所引导的观光之旅，
为处在极其不同水平上的学习者而设计。

\seclbl{练习}{s:process-exercises}

\exercise{创建学习者画像}{小组}{30}

以小组为单位，
创建一个四点式画像，描述你们的一位典型学习者。

\exercise{给学习目标分类}{两人一组}{10}

看看\secref{s:process-objectives}里
HTML 和 CSS 入门课程的示例学习目标，
按布鲁姆分类法给每一条分类。
和同伴的答案比较。
你们在哪里一致、在哪里分歧？

\exercise{写学习目标}{两人一组}{20}

用布鲁姆分类法，
为你目前教或打算教的某个东西写一个或多个学习目标。
和同伴一起
批评并改进这些目标。
每一个是否都有可验证的动词，
并清楚规定了可接受表现的标准？

\exercise{再写一些学习目标}{两人一组}{20}

用芬克分类法，
为你目前教或打算教的某个东西写一个或多个学习目标。
和同伴一起
批评并改进这些目标。

\exercise{帮我自己做成}{小组}{15}

教育理论家列夫·维果茨基提出了\gref{g:zpd}{最近发展区}（ZPD）的概念，
它包含人们目前还不能独立解决、
但在导师帮助下可以解决的一类问题。
这些正是接下来最值得着手的问题，
因为它们够不着、却又能达到。

以小组为单位，
选一个你们已经做出的学习者画像，
描述两三个处在该学习者最近发展区内的问题。

\exercise{用减复杂度来搭建课程}{个人}{20}

搭建编程课的一种办法，
是先写出你想让学习者在最后完成的程序，
然后去掉你最想让他们写的那部分最复杂的内容，
把它变成最后一道练习。
接着再去掉次复杂、也想让他们写的那部分，
把它变成倒数第二道练习，
依此类推。
把练习都抽走之后剩下的东西，
比如加载库或读取数据，
就成了你给他们的起始代码。

拿一个你想让学习者能够创建的程序或网页，
往回推，把它拆成好消化的若干部分。
一共有多少部分？
每一部分引入了什么关键想法？

\exercise{无关紧要的怪癖}{个人}{15}

Betsy Leondar-Wright 造了
「\hreffoot{http://www.classmatters.org/2006\_07/its-not-them.php}{无关紧要的怪癖}（inessential weirdness）」
这个说法，用来描述一些小团体做的事，
这些事其实并非必要，
却把还不是该团体成员的人排斥在外。
Sumana Harihareswara 后来把这一观念
用在一场关于
\hreffoot{https://www.harihareswara.net/sumana/2016/05/21/0}{开源软件里无关紧要的怪癖}的演讲里，
其中就包括使用名字晦涩的命令行工具之类的事。
花几分钟读读这些文章，
然后列一份清单，写出你认为学习者在你第一次教他们时可能遇到的无关紧要的怪癖。
其中有多少你能避免？

\exercise{PETE}{个人}{15}
\index{PETE（课程模式）}
\index{课程模式!PETE}

对编程课很管用的一种模式是 PETE：
引入\ \textbf{问}题（Problem）,
过一遍一个\textbf{例}子（Example）、
讲解\textbf{理}论（Theory），
然后再\textbf{阐}发第二个例子（Elaborate），
好让学习者看出哪些是各案例特有的、哪些适用于所有案例。
挑一件你最近教过或被教过的东西，
按这四个步骤为它设计一节短课。

\exercise{PRIMM}{个人}{15}
\index{PRIMM（课程模式）}
\index{课程模式!PRIMM}

另一种课程模式是 PRIMM~\cite{Sent2019}：
\textbf{预}测（Predict）程序的行为或输出、
\textbf{运}行它（Run）看它实际做了什么、
\textbf{探}究它为什么那样做（Investigate）——比如用调试器逐步执行或画出控制流、
\textbf{改}动它（Modify）（或其输入），
然后\textbf{重}新造一个类似的（Make）。
挑一件你最近教过或被教过的东西，
按这五个步骤为它设计一节短课。

\exercise{具体—表征—抽象}{两人一组}{15}
\index{CRA（课程模式）}
\index{课程模式!CRA}

\hreffoot{https://makingeducationfun.wordpress.com/2012/04/29/concrete-representational-abstract-cra/}{具体—表征—抽象}
（Concrete-Representational-Abstract，CRA）
是一种引入新想法的方法，
主要用在年龄较小的学习者身上：
亲手摆弄一个\textbf{具}体物体（Concrete）、
用图像\textbf{表}征这个物体（Represent）、
然后使用数字、符号或其他\textbf{抽}象物（Abstraction）
来做同样的操作。

\begin{enumerate}

\item
  把 2、7、5、10、6 这几个数字分别写在一张便利贴上。

\item
  依次看每一个，模拟一个寻找最大值的循环（具体）。

\item
  画出你所用过程的示意图，
  给每一步加上标注（表征）。

\item
  写出别人能照着走同样步骤的说明（抽象）。
\end{enumerate}

把你的表征材料和抽象材料与同伴的比较。

\exercise{评估一个课程仓库}{小组}{10}

\cite{Leak2017}探讨了计算机科学老师为什么不使用课程分享网站，
并提出了让它们更有吸引力的建议：

\begin{enumerate}

\item
  落地页应该让访客能标明自己的背景和访问该网站的兴趣所在。
  网站应该问两个问题：
  「你目前的角色是什么？」以及
  「你对哪门课程和哪个年级感兴趣？」

\item
  网站应该在完整课程的语境中展示所有学习资源，
  好让潜在用户能理解它们预期使用的语境。

\item
  很多老师担心，如果自己在网站的讨论区发帖，
  自己的（知识欠缺）会被同行评判。
  因此这些讨论区应该允许匿名发帖。

\end{enumerate}

以小组为单位，
讨论这三个特点是否足以说服你去使用一个课程分享网站；
如果不够，
还需要什么。

\section*{回顾}

\figpdfhere{figures/conceptmap-personas.pdf}{概念：学习者画像}{f:process-personas}{}
